\section*{Devoir maison }

On s'intéresse à la fonction max qui prend en paramètre un tableau de valeurs numériques non vide et qui renvoie la valeur maximum du tableau ainsi que la position à laquelle il se trouve.

\begin{lstlisting}[language=Python]
def recherche_max(tab):
    vmax, imax = 0, None
    for i in range(len(tab)-1):
        if tab[i] > vmax:
            vmax = tab[i]
            imax = i
    return vmax, imax
\end{lstlisting}

\begin{Exercise}[title={Trouver un maximum}]

\Question Quel type de parcours est utilisé à la ligne 3 de la fonction ?\\
\Question Que représente la variable imax ? \\
\Question Écrire l'appel de la fonction recherche\_max() qui prendra en paramètre la liste [9, 100, 8, 48, 73]\\
\Question Compléter le tableau d'avancement des variables pour un appel de la fonction recherche\_max() qui prendra en paramètre la liste [9, 100, 8, 48, 73].\\
 \begin{center}
\large
\begin{tabular}{c|c|c|c|c|c|}
\textbf{Tour} & \textbf{i} & \textbf{tab[i]} & \textbf{tab[i] > vmax} & \textbf{vmax} & \textbf{imax } \\
\hline
0 &  &  &  & &\\
1 &  &  &  & &\\
2 &  &  &  & &\\
3 &  &  &  & &\\
4 &  &  &  & &
\end{tabular}
\end{center}

\Question Quelles sont les valeurs de retours de la fonction recherche\_max pour les appels suivants : \\
\begin{itemize}
 \item recherche\_max([10, 4, 8, 88, 10])\\ \emptybox{1cm}
 \item recherche\_max([-10, -2, -10]) \\ \emptybox{1cm}
\end{itemize}
\Question La fonction proposée est-elle correcte ? Si non, justifiez en vous aidant de la question précédente\\[0.2cm]
\emptybox{2cm}
\Question Réécrivez la fonction en corrigeant le problème\\[0.2cm]
\emptybox{3cm}

\end{Exercise}